\begin{helper}
Let \(V\) be finite, let \(\nu\) be a finite-on-the-fiber measure on a space
\(\Omega\), let \(F\subseteq\Omega\) be a fiber event, and let
\(p:\Omega\to\mathbb R^V\) be a priority-coordinate map.  Suppose that the
fiber priority coordinates have the lower-rectangle law with mass \(M\):
for every \(t=(t_v)_{v\in V}\in[0,1]^V\),
\[
\nu\{\omega:F(\omega),\ 0\le p(\omega)_v\le t_v\text{ for every }v\}
=M\operatorname{ofReal}\Bigl(\prod_{v\in V}t_v\Bigr),
\]
where \(M=\nu\{\omega:F(\omega)\}\) and this fiber mass is finite.  Let
\(\gamma>0\), let \(a,b,c\in V\) be pairwise distinct, and let
\(S_x,S_y,S_z\subseteq V\) be pairwise disjoint and disjoint from
\(\{a,b,c\}\).  Then
\[
\begin{aligned}
&\nu\bigl(\{\omega:F(\omega),\ p_a(\omega)\le p_b(\omega)\le p_c(\omega),\\
&\qquad\qquad\quad
  p_w(\omega)<p_a(\omega)\ (w\in S_x),\
  p_w(\omega)<p_b(\omega)\ (w\in S_y),\
  p_w(\omega)<p_c(\omega)\ (w\in S_z)\}\\
&\qquad\qquad\cap\{\omega:p_v(\omega)\in[0,1]\text{ for every }v\in V\}\bigr)
\end{aligned}
\]
is at most
\[
M\operatorname{ofReal}\left(
\frac{1}{\gamma^3}
\int_{0\le z\le y\le x\le\gamma}
\left(1-\frac{x}{\gamma}\right)^{|S_x|}
\left(1-\frac{y}{\gamma}\right)^{|S_y|}
\left(1-\frac{z}{\gamma}\right)^{|S_z|}\,dx\,dy\,dz\right).
\]
\end{helper}
\begin{proof}
Restrict \(\nu\) to the fiber \(F\) and push it forward by the coordinate map
\(p\).  The lower-rectangle hypothesis says that this pushed-forward finite
measure agrees on every anchored rectangle
\(\prod_v[0,t_v]\), \(0\le t_v\le1\), with \(M\) times the product Lebesgue
probability measure on \([0,1]^V\).  These anchored rectangles form a
pi-system generating the Borel sigma-algebra of the finite cube, and the two
finite measures have the same total mass \(M\).  By the finite-measure
uniqueness theorem for generated sigma-algebras, the two measures therefore
agree on every Borel subset of the cube.

The chamber event in the statement is a Borel subset of the cube because it is
defined by finitely many coordinate inequalities.  We may consequently compute
its measure under \(M\) times product Lebesgue measure on \([0,1]^V\).  Since
the sets \(S_x,S_y,S_z\) are pairwise disjoint and avoid \(a,b,c\), Fubini's
theorem factors the outside-coordinate conditions independently once the three
coordinates \(p_a,p_b,p_c\) are fixed.  For fixed values
\(\alpha=p_a\), \(\beta=p_b\), and \(\chi=p_c\) with
\(\alpha\le\beta\le\chi\), the factors are
\(\alpha^{|S_x|}\), \(\beta^{|S_y|}\), and \(\chi^{|S_z|}\); all other
coordinates are unrestricted and contribute \(1\).  Boundary hyperplanes have
zero product measure, so the strict outside inequalities and weak ordered
coordinate inequalities are represented by the same iterated integral.

Finally apply the affine change of variables
\[
x=\gamma(1-\alpha),\qquad y=\gamma(1-\beta),\qquad z=\gamma(1-\chi).
\]
Because \(\gamma>0\), this maps the region
\(0\le\alpha\le\beta\le\chi\le1\) bijectively onto
\(0\le z\le y\le x\le\gamma\) with Jacobian factor \(\gamma^{-3}\), and it
turns the three outside-coordinate factors into
\((1-x/\gamma)^{|S_x|}\), \((1-y/\gamma)^{|S_y|}\), and
\((1-z/\gamma)^{|S_z|}\).  Multiplying by the fiber mass \(M\) gives exactly
the displayed bound.
\end{proof}
